% ============================================================================
%  3.4 MODELO DE DATOS
% ============================================================================

\section{Modelo de datos}

\subsection{Tabla \texttt{event\_registrations}}

Se cre\'o una nueva tabla en PostgreSQL mediante la migraci\'on
\texttt{011\_event\_registrations.sql}:

\begin{verbatim}
CREATE TABLE public.event_registrations (
  id UUID PRIMARY KEY DEFAULT gen_random_uuid(),
  user_id UUID NOT NULL REFERENCES public.profiles(id)
    ON DELETE CASCADE,
  event_id UUID NOT NULL REFERENCES public.events(id)
    ON DELETE CASCADE,
  status TEXT NOT NULL DEFAULT 'confirmed'
    CHECK (status IN ('confirmed', 'cancelled', 'waitlist')),
  registered_at TIMESTAMPTZ NOT NULL DEFAULT NOW(),
  UNIQUE (user_id, event_id)
);
\end{verbatim}

\subsection{Campos de la tabla}

\begin{table}[H]
\centering
\begin{tabular}{@{}lllp{5cm}@{}}
\toprule
\textbf{Campo} & \textbf{Tipo} & \textbf{Restricci\'on} &
  \textbf{Descripci\'on} \\
\midrule
id & UUID & PK, auto & Identificador \'unico \\
user_id & UUID & FK $\to$ profiles & Usuario inscrito \\
event_id & UUID & FK $\to$ events & Evento al que se inscribe \\
status & TEXT & CHECK, default & Estado: confirmed, cancelled,
  waitlist \\
registered_at & TIMESTAMPTZ & default NOW() & Fecha de inscripci\'on \\
\bottomrule
\end{tabular}
\caption{Campos de la tabla \texttt{event\_registrations}}
\end{table}

\subsection{Relaciones}

\begin{figure}[H]
\centering
\begin{tikzpicture}[
  node distance=1.5cm,
  entity/.style={rectangle, draw=black, thick, minimum width=2.2cm,
    minimum height=0.8cm, align=center, font=\footnotesize,
    fill=blue!8},
  relation/.style={-Stealth, thick}
]
  \node[entity] (profiles) {profiles};
  \node[entity, right=of profiles] (reg) {event\_\\registrations};
  \node[entity, right=of reg] (events) {events};

  \draw[relation] (profiles) -- node[above, font=\scriptsize]
    {user\_id} (reg);
  \draw[relation] (reg) -- node[above, font=\scriptsize]
    {event\_id} (events);
\end{tikzpicture}
\caption{Diagrama de relaciones del m\'odulo}
\end{figure}

\subsection{Pol\'iticas RLS (Row Level Security)}

La tabla habilita RLS con las siguientes pol\'iticas:

\begin{table}[H]
\centering
\small
\begin{tabular}{@{}lp{4.5cm}p{4.5cm}@{}}
\toprule
\textbf{Pol\'itica} & \textbf{Acci\'on} & \textbf{Condici\'on} \\
\midrule
select\_own & SELECT & user\_id = auth.uid() \\
select\_admin & SELECT & Rol = admin \\
select\_gestor & SELECT & Es gestor del escenario \\
insert\_own & INSERT & user\_id = auth.uid() \\
update\_own & UPDATE & user\_id = auth.uid() \\
delete\_own\_or\_admin & DELETE & user\_id = auth.uid() o admin \\
\bottomrule
\end{tabular}
\caption{Pol\'iticas RLS de la tabla \texttt{event\_registrations}}
\end{table}

\subsection{Vista \texttt{event\_registration\_counts}}

Se cre\'o una vista para obtener el conteo de inscripciones por evento:

\begin{verbatim}
CREATE OR REPLACE VIEW
  public.event_registration_counts AS
SELECT
  event_id,
  COUNT(*) FILTER (WHERE status = 'confirmed')
    AS confirmed_count,
  COUNT(*) FILTER (WHERE status = 'waitlist')
    AS waitlist_count
FROM public.event_registrations
GROUP BY event_id;
\end{verbatim}
